\section{Medidas de dispersión}

\subsection{Dispersión o variación}
Si bien las medidas de tendencia central nos dicen alrededor de qué valores se concentra un arreglo de datos, las \emph{medidas de dispersión} nos dan una idea de qué tan alejados están los datos entre sí y respecto a la media.

\begin{observacion}
	Dos conjuntos de datos pueden tener la misma media pero muy distinta dispersión. Por ejemplo, los conjuntos $\{49, 50, 51\}$ y $\{0, 50, 100\}$ tienen la misma media ($\bar{X} = 50$), pero claramente el segundo está mucho más disperso.
\end{observacion}

A continuación veremos algunas medidas de dispersión comúnmente usadas en estadística.

\subsection{Rango}
El \emph{rango} de un conjunto de datos es la diferencia entre el mayor y el menor valor del conjunto.

\begin{ejemplo}
	El rango del conjunto $2,3,3,5,5,5,8,10,12$ es $12-2=10.$
\end{ejemplo}

\begin{observacion}
	El rango es la medida de dispersión más simple, pero tiene la desventaja de depender únicamente de los valores extremos, ignorando la distribución de los datos intermedios.
\end{observacion}


\subsection{Desviación media}
La \emph{desviación media} o \emph{desviación promedio} de un conjunto de $N$ números $X_{1},...,X_{N}$ está definida como
\begin{align}
	DM=\dfrac{\sum \abs{X_{j}-\bar{X}}}{N}
\end{align}
donde $\bar{X}$ es la media aritmética de los números y $\abs{\cdot}$ denota el valor absoluto.

\begin{ejemplo}
	Encuentre la desviación media de la lista $2,3,6,8,11.$
\end{ejemplo}

\begin{solucion}
	Primero calculamos la media:
	\begin{align}
		\bar{X} = \frac{2+3+6+8+11}{5} = \frac{30}{5} = 6
	\end{align}
	
	Luego calculamos las desviaciones absolutas respecto a la media:
	\begin{align}
		DM &= \frac{|2-6| + |3-6| + |6-6| + |8-6| + |11-6|}{5} \\
		&= \frac{4 + 3 + 0 + 2 + 5}{5} = \frac{14}{5} = 2.8
	\end{align}
\end{solucion}


\subsection{Desviación estándar}
La \emph{desviación estándar} de un conjunto de $N$ números $X_{1},...,X_{N}$ se denota como $s$ y está definida por
\begin{align}
	s=\sqrt{\dfrac{\sum\left( X_{j}-\bar{X} \right)^{2}}{N}}=\sqrt{\sum\dfrac{x_{j}^{2}}{N}}
\end{align} donde $x_{j}:=X_{j}-\bar{X}.$

\begin{observacion}
	La desviación estándar mide qué tan dispersos están los datos respecto a la media. Una desviación estándar pequeña indica que los datos están concentrados cerca de la media, mientras que una desviación estándar grande indica que los datos están muy dispersos.
\end{observacion}

Si $X_{1},...,X_{N}$ se presentan con frecuencias $f_{1},...,f_{N}$ respectivamente, la desviación estándar se puede expresar como
\begin{align}
	s=\sqrt{\dfrac{\sum f_{j}\left( X_{j}-\bar{X} \right)^{2}}{N}}=\sqrt{\dfrac{\sum f_{j}x_{j}^{2}}{N}}
\end{align}

\begin{ejemplo}
	Calcule la desviación estándar del conjunto $2, 4, 4, 4, 5, 5, 7, 9.$
\end{ejemplo}

\begin{solucion}
	Primero calculamos la media:
	\begin{align}
		\bar{X} = \frac{2+4+4+4+5+5+7+9}{8} = \frac{40}{8} = 5
	\end{align}
	
	Luego calculamos las desviaciones al cuadrado:
	\begin{align}
		(2-5)^2 &= 9 \\
		(4-5)^2 &= 1 \quad (\text{tres veces}) \\
		(5-5)^2 &= 0 \quad (\text{dos veces}) \\
		(7-5)^2 &= 4 \\
		(9-5)^2 &= 16
	\end{align}
	
	Por lo tanto:
	\begin{align}
		s &= \sqrt{\frac{9 + 3(1) + 2(0) + 4 + 16}{8}} = \sqrt{\frac{32}{8}} = \sqrt{4} = 2
	\end{align}
\end{solucion}


\subsection{Varianza}
La \emph{varianza} de un conjunto de números se define como el cuadrado $s^{2}$ de la desviación estándar $s$.

\begin{observacion}
	En estadística es importante distinguir entre la desviación estándar de una \emph{población} y de una \emph{muestra}. Para distinguirlas, en el primer caso utilizaremos $\sigma$ y en el segundo continuaremos usando $s.$
	
	Cuando trabajamos con una muestra y queremos estimar la varianza poblacional, utilizamos $N-1$ en lugar de $N$ como denominador:
	\begin{align}
		s^2 = \frac{\sum (X_j - \bar{X})^2}{N-1}
	\end{align}
	
	Esta corrección (llamada \emph{corrección de Bessel}) produce un estimador insesgado de la varianza poblacional. Para muestras grandes ($N > 30$), la diferencia entre usar $N$ y $N-1$ es prácticamente despreciable.
\end{observacion}


\subsection{Método abreviado}
Una forma computacionalmente más eficiente de calcular la varianza es:
\begin{align}
	s^{2} &= \overline{X^{2}} - \overline{X}^{2} \\
	&= \frac{\sum X_j^2}{N} - \left(\frac{\sum X_j}{N}\right)^2
\end{align}

donde $\overline{X^2}$ denota la media de los cuadrados y $\overline{X}^2$ denota el cuadrado de la media.

\begin{observacion}
	Esta fórmula es equivalente a la definición original pero evita calcular las desviaciones individuales. Es particularmente útil para cálculos manuales o cuando se trabaja con datos agrupados.
\end{observacion}


\subsection{Coeficiente de variación}
El \emph{coeficiente de variación} es una medida de dispersión relativa que permite comparar la variabilidad de conjuntos de datos con diferentes unidades o escalas:
\begin{align}
	CV = \frac{s}{\bar{X}} \times 100\%
\end{align}

\begin{ejemplo}
	Compare la variabilidad de las estaturas y pesos de un grupo de personas:
	\begin{itemize}
		\item Estaturas: $\bar{X} = 170$ cm, $s = 8$ cm
		\item Pesos: $\bar{X} = 70$ kg, $s = 10$ kg
	\end{itemize}
\end{ejemplo}

\begin{solucion}
	\begin{align}
		CV_{\text{estatura}} &= \frac{8}{170} \times 100\% \approx 4.71\% \\
		CV_{\text{peso}} &= \frac{10}{70} \times 100\% \approx 14.29\%
	\end{align}
	
	Aunque la desviación estándar del peso ($10$ kg) es mayor que la de la estatura ($8$ cm) en valor absoluto, el coeficiente de variación nos muestra que el peso tiene mucha mayor variabilidad relativa ($14.29\%$ vs $4.71\%$).
\end{solucion}


\subsection{Varianza de conjuntos combinados}
Si dos conjuntos de $N_{1}$ y $N_{2}$ datos respectivamente tienen correspondientes varianzas $s_{1}^{2}$ y $s_{2}^{2}$ pero una misma media aritmética $\bar{X},$ entonces la varianza del conjunto combinado es
\begin{align}
	s^{2}=\dfrac{N_{1}s_{1}^{2}+N_{2}s_{2}^{2}}{N_{1}+N_{2}}.
\end{align}


\subsection{Teorema de Chebyshev}

\begin{teorema}[Chebyshev]
	Para cualquier variable aleatoria $X$ con media $\mu$ y desviación estándar $\sigma$, y para cualquier $k > 1$,
	\begin{align}
		P(|X - \mu| < k\sigma) \geq 1 - \frac{1}{k^2}
	\end{align}
	
	Es decir, por lo menos $1 - \frac{1}{k^2}$ de la distribución de probabilidad de $X$ está a no más de $k$ desviaciones estándar de la media.
\end{teorema}

\begin{proof}
	Demostraremos el teorema para el caso discreto. Sea $f(x)$ la función de probabilidad de $X$.
	
	La varianza es:
	\begin{align}
		\sigma^2 = \sum_{x} (x - \mu)^2 f(x)
	\end{align}
	
	Separamos la suma en dos partes: valores dentro de $k\sigma$ de la media y valores fuera:
	\begin{align}
		\sigma^2 &= \sum_{|x-\mu| < k\sigma} (x-\mu)^2 f(x) + \sum_{|x-\mu| \geq k\sigma} (x-\mu)^2 f(x)
	\end{align}
	
	La primera suma es no negativa, así que:
	\begin{align}
		\sigma^2 \geq \sum_{|x-\mu| \geq k\sigma} (x-\mu)^2 f(x)
	\end{align}
	
	Para los términos en la segunda suma, $(x-\mu)^2 \geq k^2\sigma^2$, por lo que:
	\begin{align}
		\sigma^2 &\geq \sum_{|x-\mu| \geq k\sigma} k^2\sigma^2 f(x) = k^2\sigma^2 \sum_{|x-\mu| \geq k\sigma} f(x) \\
		&= k^2\sigma^2 P(|X-\mu| \geq k\sigma)
	\end{align}
	
	Dividiendo entre $k^2\sigma^2$:
	\begin{align}
		P(|X-\mu| \geq k\sigma) \leq \frac{1}{k^2}
	\end{align}
	
	Por lo tanto:
	\begin{align}
		P(|X-\mu| < k\sigma) = 1 - P(|X-\mu| \geq k\sigma) \geq 1 - \frac{1}{k^2}
	\end{align}
\end{proof}

\begin{observacion}
	Para $k=2$, el teorema garantiza que al menos $1 - \frac{1}{4} = 75\%$ de los datos están dentro de 2 desviaciones estándar de la media. Para $k=3$, al menos $1 - \frac{1}{9} \approx 88.9\%$ están dentro de 3 desviaciones estándar.
	
	En distribuciones normales, estos porcentajes son mucho mayores: $95.45\%$ dentro de $\pm 2\sigma$ y $99.73\%$ dentro de $\pm 3\sigma$.
\end{observacion}


\subsection{Python}

\texttt{Numpy} proporciona funciones para calcular todas las medidas de dispersión:

\begin{lstlisting}[language=Python]
import numpy as np

data = np.array([2, 4, 4, 4, 5, 5, 7, 9])

# Rango
rango = np.ptp(data)
print(f"Rango: {rango}")  # 7

# Desviación media
desv_media = np.mean(np.abs(data - np.mean(data)))
print(f"Desviación media: {desv_media}")  # 1.5

# Desviación estándar (poblacional, ddof=0)
std_pop = np.std(data, ddof=0)
print(f"Desviación estándar poblacional: {std_pop}")  # 2.0

# Desviación estándar (muestral, ddof=1)
std_muestra = np.std(data, ddof=1)
print(f"Desviación estándar muestral: {std_muestra}")  # 2.138

# Varianza
var_pop = np.var(data, ddof=0)
print(f"Varianza poblacional: {var_pop}")  # 4.0

var_muestra = np.var(data, ddof=1)
print(f"Varianza muestral: {var_muestra}")  # 4.571

# Coeficiente de variación
cv = (std_pop / np.mean(data)) * 100
print(f"Coeficiente de variación: {cv:.2f}%")  # 40.00%
\end{lstlisting}

\begin{observacion}
	El parámetro \texttt{ddof} (Delta Degrees of Freedom) controla el denominador: \texttt{ddof=0} usa $N$ (poblacional) y \texttt{ddof=1} usa $N-1$ (muestral).
\end{observacion}
